\documentclass[11pt,a4paper]{article}
\usepackage[margin=22mm]{geometry}
\usepackage{kotex}
\usepackage{amsmath,amssymb,mathtools,bm}
\usepackage{booktabs,longtable,array,tabularx}
\usepackage{enumitem}
\usepackage{hyperref}
\usepackage{amsthm}
\usepackage{mdframed}
\usepackage{fancyhdr}
\usepackage{lastpage}
\usepackage{setspace}
\setstretch{1.13}
\setlength{\parskip}{0.4em}
\setlength{\parindent}{0pt}
\setlength{\headheight}{14pt}
\hypersetup{colorlinks=true, linkcolor=blue, urlcolor=blue, citecolor=blue,
  pdftitle={Graphite ICA tail thermodynamic extension: reversible heat and dVdT}}
\pagestyle{fancy}\fancyhf{}
\lhead{Graphite ICA tail — Chapter 2}\rhead{\thepage/\pageref{LastPage}}
\renewcommand{\headrulewidth}{0.4pt}

\newtheorem*{groundingbox}{Grounding}
\newtheorem*{boundbox}{Validity bound}
\newtheorem*{warningbox}{Warning}

\newcommand{\dd}{\mathrm{d}}
\newcommand{\eff}{\mathrm{eff}}
\newcommand{\eqs}{\mathrm{eq}}
\newcommand{\dyn}{\mathrm{dyn}}
\newcommand{\app}{\mathrm{app}}
\newcommand{\rev}{\mathrm{rev}}
\newcommand{\irr}{\mathrm{irr}}
\newcommand{\bg}{\mathrm{bg}}
\newcommand{\cell}{\mathrm{cell}}
\newcommand{\ext}{\mathrm{ext}}
\newcommand{\tail}{\mathrm{tail}}
\newcommand{\pol}{\mathrm{pol}}
\newcommand{\lag}{\mathrm{lag}}
\newcommand{\thm}{\mathrm{th}}
\newcommand{\calH}{\mathcal H}

\title{\textbf{리튬이온전지 graphite anode 의 ICA(\,$dQ/dV$\,) tail 거동:\\
reversible heat 와 \(dV/dT\) 로의 thermodynamic extension}\\[0.4em]
\large Chapter 2}
\author{Project\_Anode\_Fit}
\date{}

\begin{document}
\maketitle

\begin{center}\begin{minipage}{0.93\textwidth}
\small \textbf{Abstract.}
Chapter 1 은 graphite anode ICA peak 뒤의 post-peak tail 을 activation barrier,
electrode-potential assistance, relaxation-length spectrum \(A_L(\lambda;T,\psi)\) 로
해석했다. Chapter 2 는 그 식을 reversible heat 와 \(dV/dT\) 로 확장한다. 핵심은
\textbf{equilibrium entropy coefficient} 와 \textbf{dynamic apparent temperature derivative} 를
분리하는 것이다. Equilibrium/quasi-static limit 에서는 charge conservation 을 fixed \(q\) 에서
temperature 로 미분하여
\[
h_n^{\eqs}(q,T)
=
\left(\frac{\partial V_{n,\eqs}}{\partial T}\right)_q
=
-\frac{
\left(\partial Q_{\bg}/\partial T\right)_V
+Q_p\left(\partial\Theta_{\eqs}/\partial T\right)_V}
{C_{\bg}(V_{n,\eqs},T)+Q_p\left(\partial\Theta_{\eqs}/\partial V_n\right)_T}
\]
를 얻는다. 이 \(h_n^{\eqs}\) 가 reversible heat 의 entropy coefficient 이다. 반면
galvanostatic ICA 중 관측되는 \((\partial V_{n,\app}/\partial T)_q\) 는 equilibrium term,
kinetic-lag/history term, tail-spectrum redistribution, polarization term 이 섞인
dynamic apparent derivative 이다. 따라서 Chapter 1 의 tail kernel 은 measured thermal
response 를 설명하는 데 쓰일 수 있지만, 그 자체가 곧 reversible heat 는 아니다. 이 장의
목표는 이 경계를 수식으로 고정하여 Chapter 3 의 reaction kinetics, Chapter 4 의 heat coupling,
Chapter 5 의 charge/discharge hysteresis 로 넘어갈 수 있는 thermodynamic foundation 을 만드는 것이다.
\end{minipage}\end{center}

\tableofcontents
\newpage

% ====================================================================
\section{Target system 과 Chapter 2 의 central separation}\label{sec:target}
% ====================================================================

본 문건의 대상은 리튬이온전지(LIB) graphite anode 의 incremental capacity analysis (ICA),
즉 \(dQ/dV\) 곡선이다. 관심 구간은 staging transition 으로 생기는 peak 자체보다
peak 이후의 post-peak tail 이다. 실험적으로 temperature 가 낮을수록 tail 이 길게 늘어지고,
temperature 가 높을수록 tail 이 상대적으로 빠르게 끝난다. Chapter 1 은 이 현상을
smooth activation barrier 와 relaxation-length spectrum 으로 설명했다. Barrier 는 어떤 지점에서
상태를 갑자기 \(0\to1\) 로 바꾸는 threshold 가 아니라, local transition mode 가 equilibrium target 을
따라가는 rate constant 를 정하는 free-energy quantity 이다.

Chapter 2 에서 추가되는 질문은 다음이다.

\begin{quote}
Chapter 1 의 ICA-tail 식을 이용하여 reversible heat, entropy coefficient, 그리고
\(dV/dT\) 를 어떻게 해석할 수 있는가?
\end{quote}

이 질문에 답하려면 가장 먼저 두 quantity 를 분리해야 한다.

\[
\text{equilibrium entropy coefficient}
\quad\ne\quad
\text{dynamic apparent temperature derivative}.
\]

Equilibrium entropy coefficient 는 충분히 relaxed 된 electrode 의 equilibrium potential
\(V_{n,\eqs}\) 에 대해
\[
h_n^{\eqs}(q,T)
\equiv
\left(\frac{\partial V_{n,\eqs}}{\partial T}\right)_q
\]
로 정의된다. 이 quantity 는 reversible heat 와 직접 연결된다. 반면 galvanostatic ICA 중
측정되는 apparent voltage \(V_{n,\app}\) 의 temperature derivative 는
\[
h_n^{\app}(q,T;\calH)
\equiv
\left(\frac{\partial V_{n,\app}}{\partial T}\right)_{q,\calH}
\]
처럼 protocol/history \(\calH\) 에 의존한다. 이 값은 equilibrium thermodynamics 뿐 아니라
kinetic lag, spectrum redistribution, ohmic/contact polarization, transport polarization 을 함께 포함한다.

\begin{warningbox}
Raw apparent voltage 에서 얻은 \(dV/dT\) 를 곧바로 reversible heat 의 entropy coefficient 로
사용하면 논리 오류가 생긴다. Chapter 2 의 핵심 deliverable 은 이 오류를 막는 분리식이다.
\end{warningbox}

% ====================================================================
\section{Conventions}\label{sec:conventions}
% ====================================================================

\subsection{Coordinate 와 state variable}\label{sec:coordinate}

Chapter 2 는 Chapter 1 의 convention 을 유지한다.

\begin{longtable}{@{}p{0.22\textwidth}p{0.70\textwidth}@{}}
\toprule
Symbol & Meaning \\
\midrule
\endhead
\(q\) & normalized external charge coordinate. \(q\) 가 이 문건의 primary coordinate 이다. \\
\(Q_{\ext}\) & dimensional external charge, \(Q_{\ext}=Q_{\cell}q\). \\
\(Q_{\cell}\) & chosen cell/electrode capacity scale. \\
\(V_n\) & internal graphite-anode potential coordinate in the theory. \\
\(V_{n,\eqs}\) & equilibrium or quasi-static graphite-anode potential. \\
\(V_{n,\app}\) & measured/apparent graphite-anode potential after polarization is included. \\
\(\theta_j\) & local transition progress of mode \(j\). \\
\(\theta_{\eqs,j}\) & equilibrium target of mode \(j\). \\
\(\Theta\) & summed transition inventory entering charge conservation. \\
\(\Theta_{\eqs}\) & equilibrium transition inventory. \\
\(r_j\) & lag variable, \(r_j=\theta_{\eqs,j}-\theta_j\). \\
\(L_q\) & normalized relaxation length along \(q\). \\
\(A_L(\lambda;T,\psi)\) & relaxation-length spectrum. \\
\(\psi\) & electrode-potential assistance coordinate. \\
\bottomrule
\end{longtable}

No separate transformed progress coordinate is introduced. If a future fitting layer wants such a coordinate,
that transform must be explicitly derived from \(q\), \(V_n\), or \(\theta_j\), not silently substituted.

\subsection{Current, heat sign, and thermal quantities}\label{sec:signs}

Thermal equations are sign-sensitive, and the sign convention differs across battery literature.
This manuscript therefore carries the current sign explicitly.

\begin{longtable}{@{}p{0.22\textwidth}p{0.70\textwidth}@{}}
\toprule
Symbol & Meaning \\
\midrule
\endhead
\(I_s\) & signed current in the chosen energy-balance convention. Reversing the current convention changes the sign of \(I_s\), not the thermodynamic derivation. \\
\(|I|\) & current magnitude. Chapter 1 relaxation length uses \(|I|\). \\
\(\dot Q_{\rev,n}\) & reversible heat rate assigned to graphite anode; positive means heat generation into the cell body under the stated sign convention. \\
\(\dot Q_{\irr,n}\) & irreversible heat rate assigned to graphite anode; positive definite in magnitude form. \\
\(h_n^{\eqs}\) & equilibrium entropy coefficient, \((\partial V_{n,\eqs}/\partial T)_q\). \\
\(h_n^{\dyn}\) & dynamic internal temperature derivative, before apparent-potential polarization is added. \\
\(h_n^{\app}\) & apparent measured derivative, \((\partial V_{n,\app}/\partial T)_{q,\calH}\). \\
\(C_{\thm}\) & lumped thermal capacitance for later heat balance. \\
\(H_{\thm}\) & lumped heat-transfer coefficient to ambient. \\
\bottomrule
\end{longtable}

To avoid hiding a sign choice inside the equations, reversible heat is written as
\[
\dot Q_{\rev,n}=-I_sT h_n^{\eqs}.
\]
This is the common Bernardi/Newman-style sign form when \(I_s\) is the current sign used in the
corresponding electrochemical energy balance \cite{bernardi1985,newman2004}. If another convention is used,
the same physics is obtained by changing the sign definition of \(I_s\). The magnitude of the reversible heat
scale is
\[
|\dot Q_{\rev,n}|=|I|T|h_n^{\eqs}|.
\]

% ====================================================================
\section{Chapter 1 handoff}\label{sec:handoff}
% ====================================================================

Chapter 1 starts from charge conservation:
\begin{equation}
Q_{\cell}q
=
Q_{\bg}(V_n,T)+Q_p\Theta(V_n,q,T).
\label{eq:dynamic_charge_balance}
\end{equation}
Here \(Q_{\bg}\) is the smooth background charge storage, and \(Q_p\Theta\) is the charge carried by
graphite transition inventory. \(Q_p\) is a capacity scale for the transition inventory.

At fixed temperature and along a dynamic galvanostatic path, differentiating Eq.~\eqref{eq:dynamic_charge_balance}
with respect to \(q\) gives
\begin{equation}
\frac{\dd V_n}{\dd q}
=
\frac{Q_{\cell}-Q_p\,\dd\Theta/\dd q}{C_{\bg}(V_n,T)},
\qquad
C_{\bg}(V_n,T)\equiv
\left(\frac{\partial Q_{\bg}}{\partial V_n}\right)_T.
\label{eq:ch1_dvdq}
\end{equation}
This relation is the bridge from transition kinetics to ICA shape.

Chapter 1 then assigns each local transition mode a smooth effective barrier
\begin{equation}
\Delta G_{\eff,j}(V_n,T,\psi)
=
\Delta G_{a,j}(T)-\chi_j A_j(V_n,\psi),
\label{eq:effective_barrier}
\end{equation}
where \(A_j\) is a smooth assistance function. The rate constant may be written in Eyring form,
\begin{equation}
k_j(G,T,\psi)
=
k_{0,j}(T)
\exp\!\left[-\frac{G-W_{\psi,j}}{RT}\right].
\label{eq:rate}
\end{equation}
In the spectral notation, \(G\) is the distributed bare barrier coordinate and \(W_{\psi,j}\) is the smooth
potential-assistance work term. Thus \(G-W_{\psi,j}\) is the effective barrier entering the exponential rate.
Under galvanostatic forcing, this creates a normalized relaxation length
\begin{equation}
L_q(G,T,\psi)
=
\frac{|I|}{Q_{\cell}k_{0,j}(T)}
\exp\!\left[\frac{G-W_{\psi,j}}{RT}\right].
\label{eq:Lq}
\end{equation}

The observed post-peak tail is not one exponential length. It is a spectrum integral:
\begin{equation}
\frac{\dd\Theta_{\tail}}{\dd q}(q;T,\psi)
\simeq
\int_0^\infty
A_L(\lambda;T,\psi)\lambda^{-1}
\exp[-(q-q_a)/\lambda]\,\dd\lambda,
\qquad q>q_a.
\label{eq:tail_kernel}
\end{equation}
Low temperature increases the statistical weight of long \(L_q\) modes because the exponential mapping from
barrier to relaxation length becomes sharper. Electrode-potential assistance shifts modes toward shorter
\(L_q\) by lowering the effective barrier smoothly.

Chapter 2 does not alter this mechanism. Instead, it asks how the same thermodynamic structure should enter
temperature derivatives and heat.

% ====================================================================
\section{Equilibrium \(dV/dT\) from charge conservation}\label{sec:eq_dvdt}
% ====================================================================

\subsection{Why the equilibrium limit must be written separately}\label{sec:eq_need}

Reversible heat is a thermodynamic quantity. It belongs to a reversible or quasi-static path in which the
electrode state is relaxed to its equilibrium target at each \(q\) and \(T\). Therefore the starting point is
not Eq.~\eqref{eq:dynamic_charge_balance} with history-dependent \(\Theta\), but its equilibrium version:
\begin{equation}
Q_{\cell}q
=
Q_{\bg}(V_{n,\eqs},T)
+Q_p\Theta_{\eqs}(V_{n,\eqs},T).
\label{eq:eq_charge_balance}
\end{equation}
The subscript \(\eqs\) means that \(\Theta_{\eqs}\) is computed from equilibrium targets
\(\theta_{\eqs,j}\). No kinetic lag variable \(r_j\) appears in Eq.~\eqref{eq:eq_charge_balance}.

\subsection{Step-by-step fixed-\(q\) differentiation}\label{sec:differentiate}

Define
\begin{equation}
F(V,T;q)
=
Q_{\bg}(V,T)+Q_p\Theta_{\eqs}(V,T)-Q_{\cell}q.
\label{eq:F_definition}
\end{equation}
The equilibrium potential \(V_{n,\eqs}(q,T)\) is the implicit solution of
\begin{equation}
F(V_{n,\eqs}(q,T),T;q)=0.
\label{eq:F_zero}
\end{equation}

Now hold \(q\) fixed and differentiate Eq.~\eqref{eq:F_zero} with respect to \(T\). Since \(V_{n,\eqs}\) is
itself a function of \(T\), the chain rule gives
\begin{equation}
0
=
\left(\frac{\partial F}{\partial V}\right)_T
\left(\frac{\partial V_{n,\eqs}}{\partial T}\right)_q
+
\left(\frac{\partial F}{\partial T}\right)_V.
\label{eq:implicit_step}
\end{equation}
The two partial derivatives are
\begin{align}
\left(\frac{\partial F}{\partial V}\right)_T
&=
\left(\frac{\partial Q_{\bg}}{\partial V}\right)_T
+Q_p\left(\frac{\partial\Theta_{\eqs}}{\partial V}\right)_T
=
C_{\bg}(V_{n,\eqs},T)
+Q_p\left(\frac{\partial\Theta_{\eqs}}{\partial V}\right)_T,
\label{eq:FV}\\
\left(\frac{\partial F}{\partial T}\right)_V
&=
\left(\frac{\partial Q_{\bg}}{\partial T}\right)_V
+Q_p\left(\frac{\partial\Theta_{\eqs}}{\partial T}\right)_V.
\label{eq:FT}
\end{align}
Substituting Eqs.~\eqref{eq:FV} and \eqref{eq:FT} into Eq.~\eqref{eq:implicit_step} gives
\begin{equation}
h_n^{\eqs}(q,T)
\equiv
\left(\frac{\partial V_{n,\eqs}}{\partial T}\right)_q
=
-\frac{
\left(\partial Q_{\bg}/\partial T\right)_V
+Q_p\left(\partial\Theta_{\eqs}/\partial T\right)_V}
{C_{\bg}(V_{n,\eqs},T)
+Q_p\left(\partial\Theta_{\eqs}/\partial V\right)_T}.
\label{eq:h_eq}
\end{equation}

Eq.~\eqref{eq:h_eq} is the central Chapter 2 equilibrium equation. Its denominator has units of
\(\mathrm{C/V}\), because \(C_{\bg}\) is differential charge storage and
\(Q_p(\partial\Theta_{\eqs}/\partial V)_T\) is transition differential capacity. Its numerator has units of
\(\mathrm{C/K}\). Therefore the ratio has units of \(\mathrm{V/K}\), as required for \(dV/dT\).

\begin{groundingbox}
Eq.~\eqref{eq:h_eq} is just implicit differentiation of charge conservation. No fitting assumption is used.
The only physical assumption is that the state is equilibrium or quasi-static, so that
\(\Theta\) can be replaced by \(\Theta_{\eqs}\).
\end{groundingbox}

\subsection{Physical reading of the numerator and denominator}\label{sec:physical_reading}

The denominator of Eq.~\eqref{eq:h_eq},
\begin{equation}
C_{\eqs}(q,T)
\equiv
C_{\bg}(V_{n,\eqs},T)
+Q_p\left(\frac{\partial\Theta_{\eqs}}{\partial V}\right)_T,
\label{eq:Ceq}
\end{equation}
is the equilibrium differential capacity-like response. Near a staging transition, this term becomes large
because a small potential change moves a large amount of transition inventory.

The numerator,
\begin{equation}
B_{\eqs}(q,T)
\equiv
\left(\frac{\partial Q_{\bg}}{\partial T}\right)_V
+Q_p\left(\frac{\partial\Theta_{\eqs}}{\partial T}\right)_V,
\label{eq:Beq}
\end{equation}
is the charge-inventory shift caused by temperature at fixed potential. If temperature changes the equilibrium
target \(\theta_{\eqs,j}\), then holding \(q\) fixed requires the potential to shift in the opposite direction
so that total charge remains conserved. That compensation is exactly the minus sign in Eq.~\eqref{eq:h_eq}.

This explains why the entropy coefficient can have sharp features near the same graphite transition region
that creates ICA peaks. The same transition inventory appears in both denominator and numerator, but through
different partial derivatives.

% ====================================================================
\section{Thermodynamic meaning and reversible heat}\label{sec:rev_heat}
% ====================================================================

\subsection{Entropy coefficient as an operational thermodynamic quantity}\label{sec:entropy_coeff}

For a reversible electrochemical insertion reaction, electrode potential is linked to reaction free energy.
Taking a temperature derivative of that reversible potential gives an entropy coefficient. Because sign
conventions differ for anode potential, cathode potential, full-cell voltage, charge current, and discharge
current, this manuscript uses the operational definition
\begin{equation}
h_n^{\eqs}(q,T)
\equiv
\left(\frac{\partial V_{n,\eqs}}{\partial T}\right)_q.
\label{eq:h_operational}
\end{equation}
The usual thermodynamic relation can be written without hiding the convention. Let the reversible molar
free-energy change for the graphite electrode be
\begin{equation}
\Delta \bar G_n(q,T)=\sigma_VF V_{n,\eqs}(q,T),
\qquad \sigma_V\in\{+1,-1\},
\label{eq:Gbar_sign}
\end{equation}
where \(\sigma_V\) records the chosen electrode-potential orientation. Then
\begin{equation}
\Delta \bar S_n(q,T)
=
-\left(\frac{\partial \Delta \bar G_n}{\partial T}\right)_q
=
-\sigma_VF h_n^{\eqs}(q,T).
\label{eq:Sbar_sign}
\end{equation}
If the signed molar reaction rate is written as \(\dot n_s=\sigma_I I_s/F\), the reversible heat rate is
\begin{equation}
\dot Q_{\rev,n}
=
T\dot n_s\Delta \bar S_n
=
-(\sigma_I\sigma_V)I_sT h_n^{\eqs}.
\label{eq:qrev_sign_general}
\end{equation}
The convention used below absorbs \(\sigma_I\sigma_V=1\) into the definition of \(I_s\).
The connection to reversible heat is then made with the signed-current convention stated in
Section~\ref{sec:signs}:
\begin{equation}
\dot Q_{\rev,n}(q,T)
=
-I_sT h_n^{\eqs}(q,T).
\label{eq:qrev}
\end{equation}
This is the electrode-level form of the entropy heat term used in battery energy-balance theory
\cite{bernardi1985,thomas2003,newman2004}.

Eq.~\eqref{eq:qrev} should be read in two layers.

\begin{enumerate}[leftmargin=2em]
\item The thermodynamic layer is \(h_n^{\eqs}\), which is obtained from equilibrium potential.
\item The sign layer is \(I_s\), which depends on the adopted current convention.
\end{enumerate}

If another report writes reversible heat as \(+I_sT(\partial U/\partial T)\), the difference is usually a
definition of \(I_s\) or of electrode potential orientation, not a different thermodynamic principle.

\subsection{Why \(V_{n,\app}\) cannot be used directly}\label{sec:not_app}

Measured potential during current flow may be written schematically as
\begin{equation}
V_{n,\app}(q,T;\calH)
=
V_n(q,T;\calH)+\eta_n(q,T;\calH),
\label{eq:Vapp}
\end{equation}
where \(\eta_n\) collects ohmic, contact, charge-transfer, transport, and other polarization terms. Therefore
\begin{equation}
\left(\frac{\partial V_{n,\app}}{\partial T}\right)_{q,\calH}
=
\left(\frac{\partial V_n}{\partial T}\right)_{q,\calH}
+
\left(\frac{\partial\eta_n}{\partial T}\right)_{q,\calH}.
\label{eq:app_derivative}
\end{equation}
The second term is not reversible entropy. It contains resistance temperature dependence, exchange-current
temperature dependence, transport temperature dependence, and history-dependent lag. Thus
\[
h_n^{\app}=h_n^{\eqs}
\quad\text{only after equilibrium relaxation and polarization correction.}
\]

\begin{warningbox}
Apparent \(dV/dT\) is experimentally useful, but it is not automatically an entropy coefficient. Treating it
as reversible heat without removing polarization and kinetic lag double-counts dynamic effects.
\end{warningbox}

% ====================================================================
\section{Dynamic apparent \(dV/dT\)}\label{sec:dynamic_dvdt}
% ====================================================================

\subsection{History-dependent charge balance}\label{sec:history_balance}

During finite-current ICA, the transition inventory is not necessarily at equilibrium. The state depends on the
protocol history \(\calH\), including current, prior temperature, and prior potential path:
\begin{equation}
Q_{\cell}q
=
Q_{\bg}(V_n,T)+Q_p\Theta(V_n,q,T;\calH).
\label{eq:dyn_balance}
\end{equation}
The fixed-\(q\) temperature derivative follows the same implicit method as Section~\ref{sec:eq_dvdt}, but the
state is now dynamic. Define
\begin{equation}
F_{\dyn}(V,T;q,\calH)
=
Q_{\bg}(V,T)+Q_p\Theta(V,q,T;\calH)-Q_{\cell}q.
\label{eq:Fdyn}
\end{equation}
Then
\begin{equation}
h_n^{\dyn}(q,T;\calH)
\equiv
\left(\frac{\partial V_n}{\partial T}\right)_{q,\calH}
=
-\frac{
\left(\partial Q_{\bg}/\partial T\right)_V
+Q_p\left(\partial\Theta/\partial T\right)_{V,q,\calH}}
{C_{\bg}(V_n,T)+Q_p\left(\partial\Theta/\partial V\right)_{T,q,\calH}}.
\label{eq:h_dyn}
\end{equation}
This resembles Eq.~\eqref{eq:h_eq}, but the meaning is different. Eq.~\eqref{eq:h_eq} is a thermodynamic
equilibrium derivative. Eq.~\eqref{eq:h_dyn} is a protocol-dependent dynamic derivative.

\subsection{Decomposition into equilibrium, lag, and polarization pieces}\label{sec:decomp}

Write the dynamic inventory as
\begin{equation}
\Theta(V_n,q,T;\calH)
=
\Theta_{\eqs}(V_n,T)-R_{\Theta}(V_n,q,T;\calH),
\label{eq:Theta_lag}
\end{equation}
where \(R_{\Theta}\) is the total inventory lag. In a single mode, \(R_{\Theta}\) is analogous to
\(r_j=\theta_{\eqs,j}-\theta_j\); for a spectrum, it is the summed lag across modes.

Substituting Eq.~\eqref{eq:Theta_lag} into Eq.~\eqref{eq:h_dyn} shows why dynamic \(dV/dT\) differs from
equilibrium \(dV/dT\). The numerator contains
\begin{equation}
\left(\frac{\partial\Theta}{\partial T}\right)_{V,q,\calH}
=
\left(\frac{\partial\Theta_{\eqs}}{\partial T}\right)_V
-
\left(\frac{\partial R_{\Theta}}{\partial T}\right)_{V,q,\calH}.
\label{eq:Theta_T_split}
\end{equation}
The denominator contains
\begin{equation}
\left(\frac{\partial\Theta}{\partial V}\right)_{T,q,\calH}
=
\left(\frac{\partial\Theta_{\eqs}}{\partial V}\right)_T
-
\left(\frac{\partial R_{\Theta}}{\partial V}\right)_{T,q,\calH}.
\label{eq:Theta_V_split}
\end{equation}
Therefore the same graphite transition can appear in two ways:
\begin{enumerate}[leftmargin=2em]
\item as equilibrium entropy structure through \(\Theta_{\eqs}\);
\item as dynamic tail and lag structure through \(R_{\Theta}\).
\end{enumerate}
The measured apparent derivative then adds polarization:
\begin{equation}
h_n^{\app}(q,T;\calH)
=
h_n^{\dyn}(q,T;\calH)
+
\left(\frac{\partial \eta_n}{\partial T}\right)_{q,\calH}.
\label{eq:h_app}
\end{equation}

Eq.~\eqref{eq:h_app} is the practical warning equation. If a thermal experiment measures
\((\partial V_{n,\app}/\partial T)_q\) under current, it measures \(h_n^{\app}\), not necessarily
\(h_n^{\eqs}\).

% ====================================================================
\section{How the Chapter 1 tail kernel enters thermal response}\label{sec:tail_thermal}
% ====================================================================

\subsection{Tail kernel as dynamic inventory redistribution}\label{sec:tail_inventory}

In Chapter 1, the post-peak tail contribution is
\begin{equation}
g_{\tail}(q;T,\psi)
\equiv
\frac{\dd\Theta_{\tail}}{\dd q}
\simeq
\int_0^\infty
A_L(\lambda;T,\psi)\lambda^{-1}
\exp[-(q-q_a)/\lambda]\,\dd\lambda.
\label{eq:g_tail}
\end{equation}
This is a dynamic inventory-shape equation. It describes how much transition progress appears per unit
charge coordinate after the peak. Since \(A_L\) depends on \(T\) and \(\psi\), the tail shape also changes
with temperature and electrode-potential assistance.

In \(\lambda\)-space, \(\lambda\) is the integration coordinate. If the post-peak coordinate is
peak-aligned, \(s=q-q_a\), then \(s\) is held fixed and the explicit temperature dependence of
Eq.~\eqref{eq:g_tail} enters through \(A_L\):
\begin{equation}
\left(\frac{\partial g_{\tail}}{\partial T}\right)_{s,\psi}
\simeq
\int_0^\infty
\left(\frac{\partial A_L(\lambda;T,\psi)}{\partial T}\right)_{\lambda,\psi}
\lambda^{-1}
\exp[-s/\lambda]\,\dd\lambda.
\label{eq:g_tail_T_lambda}
\end{equation}
This form is clean because the kernel is written in the already transformed relaxation-length coordinate.
If the derivative is instead taken at fixed absolute \(q\), and the peak anchor \(q_a(T,\psi)\) shifts with
temperature, the chain rule adds
\begin{equation}
\left(\frac{\partial g_{\tail}}{\partial T}\right)_{q,\psi}
\simeq
\int_0^\infty
\left[
\left(\frac{\partial A_L}{\partial T}\right)_{\lambda,\psi}\lambda^{-1}
+
A_L(\lambda;T,\psi)\lambda^{-2}
\left(\frac{\partial q_a}{\partial T}\right)_{\psi}
\right]
\exp[-(q-q_a)/\lambda]\,\dd\lambda.
\label{eq:g_tail_T_absolute}
\end{equation}
Therefore a measurement must state whether tails are compared in absolute \(q\) or after peak alignment.

In the original barrier coordinate \(G\), the same derivative is less compact because the length itself depends
on \(T\):
\begin{equation}
\lambda=L_q(G,T,\psi)
=
\frac{|I|}{Q_{\cell}k_0(T)}
\exp\!\left[\frac{G-W_{\psi}}{RT}\right].
\label{eq:lambda_G}
\end{equation}
Then temperature changes both the barrier-to-length mapping and the barrier distribution. Schematically,
\begin{equation}
\frac{\partial}{\partial T}
\left[
\rho_G(G;T,\psi)
\frac{1}{L_q(G,T,\psi)}
\exp\!\left(-\frac{q-q_a}{L_q(G,T,\psi)}\right)
\right]
\label{eq:G_space_derivative}
\end{equation}
contains terms from \(\partial_T\rho_G\), \(\partial_Tk_0\), and
\(\partial_T\exp[(G-W_\psi)/(RT)]\). This is why \(\lambda\)-space is the safer manuscript form for the
thermal extension.

\subsection{Low-temperature and high-temperature tail interpretation}\label{sec:temperature_tail}

Eq.~\eqref{eq:lambda_G} shows the qualitative trend. For modes with positive effective barrier,
lower temperature increases \(L_q\) exponentially. A broader weight at large \(\lambda\) makes the post-peak
tail long. Higher temperature reduces the exponential sensitivity and shifts the observable spectrum toward
shorter relaxation lengths. Electrode-potential assistance \(W_\psi\) has a similar shortening direction:
larger assistance lowers the effective barrier and decreases \(L_q\).

The important Chapter 2 consequence is not that tail area equals reversible heat. The correct statement is:
\begin{quote}
Temperature changes the dynamic inventory lag, and that lag changes apparent \(dV/dT\) and measured thermal
response. Reversible heat, however, remains tied to equilibrium \(h_n^{\eqs}\).
\end{quote}

Thus the same transition region can create both an entropy-coefficient feature and a dynamic tail feature, but
they are not the same mathematical object.

% ====================================================================
\section{Two-layer reduced closure}\label{sec:two_layer}
% ====================================================================

\subsection{Layer E: equilibrium thermodynamic layer}\label{sec:layer_E}

Layer E is the conservative thermodynamic layer. It uses equilibrium quantities only:
\[
\Theta_{\eqs},\quad V_{n,\eqs},\quad h_n^{\eqs},\quad \dot Q_{\rev,n}.
\]
The exact implicit expression is Eq.~\eqref{eq:h_eq}. If a reduced basis is needed for interpretation or
fitting, it must be introduced as an approximation to \(h_n^{\eqs}\), not as a replacement for the definition.

Let \(\phi_j^{\eqs}(q,T)\) be normalized equilibrium transition shapes:
\begin{equation}
\phi_j^{\eqs}(q,T)
=
\frac{\left|\dd\theta_{\eqs,j}/\dd q\right|}
{\int \left|\dd\theta_{\eqs,j}/\dd q\right|\,\dd q}.
\label{eq:phi_eq}
\end{equation}
Then a reduced entropy-coefficient expansion may be written as
\begin{equation}
h_n^{\eqs}(q,T)
\approx
h_{\bg}^{\eqs}(q,T)
+
\sum_j b_j(T)\phi_j^{\eqs}(q,T).
\label{eq:h_eq_basis}
\end{equation}
Here \(b_j(T)\) has units of \(\mathrm{V/K}\). Eq.~\eqref{eq:h_eq_basis} is a basis approximation to the
thermodynamic result, and its validity must be checked against equilibrium OCV temperature-step data.

\subsection{Layer D: dynamic apparent layer}\label{sec:layer_D}

Layer D describes finite-current apparent response. It may use Chapter 1 tail kernels:
\begin{equation}
h_n^{\app}(q,T;\calH)
\approx
h_n^{\eqs}(q,T)
+
h_{\lag}(q,T;\calH)
+
h_{\pol}(q,T;\calH),
\label{eq:h_app_decomp}
\end{equation}
where
\begin{align}
h_{\lag}(q,T;\calH)
&\text{ is generated by }R_{\Theta}\text{ and the tail spectrum},\\
h_{\pol}(q,T;\calH)
&=
\left(\frac{\partial \eta_n}{\partial T}\right)_{q,\calH}.
\end{align}
A minimal tail-shape expansion is
\begin{equation}
h_{\lag}(q,T;\calH)
\approx
\sum_m c_m(T,\calH)\,
\int_0^\infty
A_{L,m}(\lambda;T,\psi)
\lambda^{-1}\exp[-(q-q_{a,m})/\lambda]\,\dd\lambda.
\label{eq:h_lag_basis}
\end{equation}
The coefficients \(c_m\) are not entropy coefficients. They convert dynamic inventory lag into apparent
temperature response under a specified protocol. This distinction is essential: Layer D is useful for fitting
measured finite-current thermal/voltage behavior, but it does not replace Layer E.

\begin{boundbox}
Layer E can be used to discuss reversible heat. Layer D can be used to discuss apparent heat-shape and
finite-current \(dV/dT\). A manuscript, patent claim, or model should not move a parameter from Layer D into
Layer E without an equilibrium reduction or an independent measurement.
\end{boundbox}

% ====================================================================
\section{Irreversible heat and double-counting boundary}\label{sec:irr_heat}
% ====================================================================

Irreversible heat is associated with finite driving force and dissipative processes. With the apparent potential
split in Eq.~\eqref{eq:Vapp}, a reduced electrode-level expression is
\begin{equation}
\dot Q_{\irr,n}
\simeq
|I|\,|\eta_n(q,T;\calH)|.
\label{eq:qirr_eta}
\end{equation}
If \(\eta_n\) is dominated by a positive effective resistance \(R_n\), then
\begin{equation}
\dot Q_{\irr,n}
\simeq
|I|^2R_n(q,T;\calH).
\label{eq:qirr_R}
\end{equation}
Eq.~\eqref{eq:qirr_R} is an approximation, not a definition. It groups several dissipative mechanisms into
one resistance-like term.

The lag variable \(R_{\Theta}\) can also reflect dissipation because a state that relaxes toward equilibrium
under finite driving force produces entropy. However, Chapter 2 does not need a separate microscopic heat
law for that lag. The conservative boundary is:
\begin{enumerate}[leftmargin=2em]
\item reversible heat uses \(h_n^{\eqs}\);
\item ohmic and overpotential losses enter \(\dot Q_{\irr,n}\);
\item kinetic-lag tail terms enter apparent response unless a later chapter derives their dissipative heat
      from a specific reaction-kinetic model.
\end{enumerate}
This avoids double counting the same finite-current effect as both reversible entropy heat and irreversible
polarization heat.

% ====================================================================
\section{Lumped thermal balance as a handoff}\label{sec:thermal_balance}
% ====================================================================

Once reversible and irreversible terms are separated, a later heat-coupled model may write
\begin{equation}
C_{\thm}\frac{\dd T}{\dd t}
=
\dot Q_{\rev,n}
+
\dot Q_{\irr,n}
+
\dot Q_{\other}
-
H_{\thm}(T-T_{\amb}).
\label{eq:lumped_heat}
\end{equation}
Here \(\dot Q_{\other}\) may include cathode heat, electrolyte heat, contact heat, side reactions, or external
heat inputs. Eq.~\eqref{eq:lumped_heat} is a bookkeeping equation for heat coupling. This chapter does not
solve it and does not use it to fit temperature.

The reason to include Eq.~\eqref{eq:lumped_heat} is architectural. Chapter 1 explains the ICA tail.
Chapter 2 identifies which part of the temperature response is equilibrium entropy and which part is dynamic
or dissipative. Chapter 4 can later insert Chapter 3 reaction kinetics into this heat balance without changing
the definitions made here.

% ====================================================================
\section{Measurement and identifiability}\label{sec:measurement}
% ====================================================================

\subsection{How to estimate \(h_n^{\eqs}\)}\label{sec:measure_eq}

The clean measurement for \(h_n^{\eqs}\) is an equilibrium or near-equilibrium temperature-step OCV method.
At a fixed \(q\), the electrode is relaxed, temperature is changed by a small amount, the electrode is relaxed
again, and the equilibrium potential shift is measured:
\begin{equation}
h_n^{\eqs}(q,T)
\approx
\frac{V_{n,\eqs}(q,T+\Delta T)-V_{n,\eqs}(q,T)}{\Delta T}.
\label{eq:measure_heq}
\end{equation}
The word ``relaxed'' is not a detail. Without sufficient relaxation, the measurement contains \(h_{\lag}\).

\subsection{What galvanostatic thermal ICA measures}\label{sec:measure_dyn}

In finite-current ICA, the measurement is closer to
\begin{equation}
\frac{V_{n,\app}(q,T_2;\calH_2)-V_{n,\app}(q,T_1;\calH_1)}{T_2-T_1}.
\label{eq:measure_happ}
\end{equation}
Unless \(\calH_1\) and \(\calH_2\) are matched and polarization is corrected, Eq.~\eqref{eq:measure_happ}
does not isolate \(h_n^{\eqs}\). It is still valuable, but its interpretation belongs to Layer D.

\subsection{Identifiability table}\label{sec:identifiability}

\begin{longtable}{@{}p{0.24\textwidth}p{0.32\textwidth}p{0.34\textwidth}@{}}
\toprule
Quantity & Best evidence & Main risk \\
\midrule
\endhead
\(h_n^{\eqs}\) & relaxed OCV temperature-step or quasi-static equilibrium map & contamination by incomplete relaxation \\
\(h_{\lag}\) & rate-dependent ICA tail change across \(T\), \(I\), and \(\psi\) & confusion with entropy coefficient \\
\(h_{\pol}\) & impedance, pulse, resistance, or overpotential correction & grouping multiple polarization mechanisms into one term \\
\(A_L(\lambda;T,\psi)\) & post-peak tail shape and its rate/temperature dependence & non-unique inversion from finite data \\
\(\dot Q_{\irr,n}\) & calorimetry plus polarization model & double counting with lag or reversible heat \\
\bottomrule
\end{longtable}

The table states a practical rule: a dynamic ICA tail feature can support the existence of a relaxation-length
spectrum, but it does not by itself identify the equilibrium entropy coefficient.

% ====================================================================
\section{Validity bounds and failure modes}\label{sec:validity}
% ====================================================================

The Chapter 2 derivation rests on the following assumptions.

\begin{enumerate}[leftmargin=2em]
\item \textbf{Local thermal uniformity.} The electrode region represented by \(T\) is treated as locally
      isothermal. Large through-thickness gradients require a spatial thermal model.
\item \textbf{Equilibrium limit for reversible heat.} Eq.~\eqref{eq:h_eq} and Eq.~\eqref{eq:qrev} require
      equilibrium or quasi-static states.
\item \textbf{Smooth assistance.} Electrode-potential assistance changes the effective barrier smoothly; it
      is not a discontinuous switch.
\item \textbf{Layer separation.} Dynamic tail terms are not called reversible heat unless an equilibrium
      reduction is shown.
\item \textbf{Finite-current history.} \(h_n^{\dyn}\) and \(h_n^{\app}\) depend on protocol history \(\calH\).
\item \textbf{Reduced basis caution.} Eq.~\eqref{eq:h_eq_basis} and Eq.~\eqref{eq:h_lag_basis} are reduced
      representations. They do not replace the definitions in Eq.~\eqref{eq:h_eq} and Eq.~\eqref{eq:h_app}.
\end{enumerate}

Failure is expected if one tries to infer \(h_n^{\eqs}\) from finite-current voltage shifts without relaxation
or polarization correction. Failure is also expected if a single Gaussian peak is forced to represent the
post-peak tail, because Chapter 1 already shows that a barrier distribution maps exponentially into a
relaxation-length spectrum.

% ====================================================================
\section{Chapter handoff}\label{sec:next}
% ====================================================================

Chapter 2 produces three outputs for later chapters:
\begin{align}
h_n^{\eqs}(q,T)
&=
-\frac{
\left(\partial Q_{\bg}/\partial T\right)_V
+Q_p\left(\partial\Theta_{\eqs}/\partial T\right)_V}
{C_{\bg}(V_{n,\eqs},T)
+Q_p\left(\partial\Theta_{\eqs}/\partial V\right)_T},
\label{eq:handoff_heq}\\
\dot Q_{\rev,n}
&=
-I_sT h_n^{\eqs},
\label{eq:handoff_qrev}\\
h_n^{\app}
&=
h_n^{\dyn}
+
\left(\partial\eta_n/\partial T\right)_{q,\calH}.
\label{eq:handoff_happ}
\end{align}

Chapter 3 can now expand \(k_j\), \(\eta_n\), and \(R_{\Theta}\) using electrochemical reaction kinetics.
Chapter 4 can insert those kinetic terms into heat generation and thermal balance. Chapter 5 can use the
same separation to interpret charge/discharge hysteresis: an equilibrium entropy feature should be symmetric
under proper equilibrium mapping, while dynamic lag and polarization terms need not be symmetric.

% ====================================================================
\section{Conclusion}\label{sec:conclusion}
% ====================================================================

Chapter 2 fixes the thermodynamic boundary required for a paper-level graphite ICA theory. The equilibrium
temperature derivative is obtained by fixed-\(q\) implicit differentiation of equilibrium charge conservation.
That derivative, \(h_n^{\eqs}\), is the entropy coefficient used for reversible heat. The finite-current
apparent derivative, \(h_n^{\app}\), is a different quantity because it contains kinetic lag, tail-spectrum
redistribution, and polarization. Chapter 1's relaxation-length spectrum remains central, but its Chapter 2
role is to explain dynamic thermal response and apparent \(dV/dT\), not to replace equilibrium thermodynamics.

\begin{thebibliography}{9}

\bibitem{bernardi1985}
D. Bernardi, E. Pawlikowski, and J. Newman,
``A General Energy Balance for Battery Systems,''
\emph{Journal of The Electrochemical Society}, 132(1), 5--12 (1985).
DOI: \href{https://doi.org/10.1149/1.2113792}{10.1149/1.2113792}.

\bibitem{thomas2003}
K. E. Thomas and J. Newman,
``Thermal Modeling of Porous Insertion Electrodes,''
\emph{Journal of The Electrochemical Society}, 150(2), A176--A192 (2003).
DOI: \href{https://doi.org/10.1149/1.1531194}{10.1149/1.1531194}.

\bibitem{newman2004}
J. Newman and K. E. Thomas-Alyea,
\emph{Electrochemical Systems}, 3rd ed.,
John Wiley \& Sons (2004).

\bibitem{degroot1984}
S. R. de Groot and P. Mazur,
\emph{Non-Equilibrium Thermodynamics},
Dover Publications (1984).

\bibitem{bard2001}
A. J. Bard and L. R. Faulkner,
\emph{Electrochemical Methods: Fundamentals and Applications}, 2nd ed.,
John Wiley \& Sons (2001).

\end{thebibliography}

\end{document}
